\section{第3課 ここは　デパートです}

\subsection{基本課文}

\begin{enumerate}
    \item ここは　デパートです。
    \item 食堂は　なな階です。
    \item あそこも　日中企画の　ビルです。
    \item かばん　売り場は　1階ですか、2階ですか。
\end{enumerate}

\begin{itemize}
    \item 甲：トイレは　どこですか。
    \item 乙：あちらです。
\end{itemize}


\begin{itemize}
    \item 甲：ここは　郵便局ですか、銀行ですか。
    \item 乙：銀行です。
\end{itemize}

\begin{itemize}
    \item 甲：これは　いくらですか。
    \item 乙：それは　５,８００円です。
    \item 甲：あれは？
    \item 乙：あれも　５,８００円です。
\end{itemize}


\subsection{语法}

\subsection{表达}

\subsection{应用课文:ホテルの周辺}


\begin{vocablist}[2]
    \vocab{デパート}
    \vocab{食堂}
    \vocab{郵便局}
    \vocab{銀行}
    \vocab{図書館}
    \vocab{マンション}
    \vocab{ホテル}
    \vocab{コンビニ}
    \vocab{喫茶店}
    \vocab{病院}
    \vocab{本屋}
    \vocab{大学}
    \vocab{レストラン}
    \vocab{ビル}
    \vocab{売り場}
    \vocab{トイレ}
    \vocab{入口}
    \vocab{事務所}
    \vocab{受付}
    \vocab{バーゲン会場}
    \vocab{エスカレーター}
    \vocab{服}
    \vocab{コート}
    \vocab{デジカメ}
    \vocab{国}
    \vocab{地図}
    \vocab{隣}
    \vocab{周辺}
    \vocab{今日}
    \vocab{水曜日}
    \vocab{木曜日}
    \vocab{ここ、そこ、あそこ}
    \vocab{こちら、そちら、あちら}
    \vocab{どこ、どちら}
    \vocab{上海}
    \vocab{東京}
    \vocab{あのう}
    \vocab{いくら}
    \vocab{お〜、〜階、〜円、〜曜日}
\end{vocablist}
